\begin{textsample}{NEG Dim 3 – 2009\_c5 – Score 1.00 – 2009\_c5/enade\_2009\_des\_2103207191.txt}  \label{ex:f3_neg_109}
Capítulo 5 Distribuição dos Conceitos 5.1 Panorama nacional da distribuição dos conceitos A Tabela 5.1 apresenta a quantidade e o percentual de cursos de DESIGN participantes do ENADE/2009 que se situaram em cada uma das faixas de conceitos, distribuídos de acordo com a Região em que se localizam os cursos das instituições de educação superior, mantenedoras dos cursos responsáveis pela oferta.
O Gráfico 5.1 apresenta a mesma informação sob a forma de diagrama de caixa para uma melhor visualização da dispersão dos conceitos obtidos.
No diagrama de caixa, a caixa propriamente dita fica limitada pelos quartis.
A barra dentro da caixa em negrito representa a mediana.
Linhas verticais unem os quartis aos valores extremos.
Se estes valores extremos se configurarem como valores atípicos, estes aparecem como pontos isolados e as linhas terminam no maior ( ou menor ) valor ainda não considerado atípico.
Eventuais valores atípicos aparecem como pequenos círculos com a Categoria Administrativa correspondente.
Valores considerados muito atípicos são representados por estrelas.
Um diagrama de caixa para poucos dados ou com dados muito concentrados pode ter alguns pontos coincidentes.
Por exemplo, se mais de metade dos cursos de uma Região receberem conceito 5, o valor extremo superior, o quartil superior e a mediana coincidirão.
Se o conjunto só contém um dado, este aparece denotado como a barra da mediana.
Se o conjunto apresenta todas as notas iguais, a notação é a mesma, somente uma barra horizontal.
Observando -se os dados da Tabela 5.1, percebe -se que, em termos nacionais, predominaram os cursos que receberam conceito 3, correspondendo respectivamente a 29,3 \% do total nacional.
Foram classificados com conceito 4, 22,7 \% dos cursos e ficaram com conceito 2, 14,7 \%.
Os conceitos máximo e mínimo foram alocados a, respectivamente, 6,7 \% e 7,3 \% dos cursos.
Ficaram sem conceito um total de 29 cursos, correspondendo a 19,3 \% do total.
Considerando -se separadamente as cinco Regiões ( ver também Gráfico 5.1 ), observa -se que, da Região Norte, 8 cursos participaram, dos quais 3 ( 37,5 \% ) obtiveram conceito 3.
Os conceitos 2 e 4 foram obtidos por 2 cursos cada ( 25,0 \% ).
Apenas 1 curso obteve o conceito mínimo e nenhum curso ficou com o conceito máximo.
Nenhum curso ficou sem conceito.
Na Região Nordeste, 15 cursos participaram, dos quais 5 obtiveram conceito 4, 3 obtiveram o conceito 3 e 2 o conceito 2.
Apenas 1 curso obteve conceito máximo e 2 cursos obtiveram conceito mínimo.
Ficaram sem conceito 2 cursos ( 13,3 \% ).
Na Região Sudeste, 69 cursos participaram, dos quais 27 obtiveram conceito 3, equivalentes a 39,1 \% do total regional.
O conceito 2 foi alcançado por 13 cursos ( 18,8 \% ), 10 cursos ( 14,5 \% ) obtiveram conceito 4 e 6 ( 8,7 \% ) cursos o conceito mínimo.
O conceito máximo foi obtido por 1 curso ( 1,4 \% ) e 12 cursos ( 17,4 \% ) ficaram sem conceito.
Na Região Sul, 16 dos 52 cursos participantes, equivalentes a 30,8 \% do total regional, classificaram -se com conceito 4.
O conceito máximo foi alcançado por 7 cursos ( 13,5 \% ), 2 cursos ( 3,8 \% ) obtiveram conceito mínimo e 10 ( 19,2 \% ) cursos ficaram com conceito 3.
Um total de 13 cursos nesta Região ( 25,0 \% ) ficaram sem conceito.
Esta Região foi a que apresentou cursos com a melhor distribuição de conceitos.
Na Região Centro-Oeste, dos 6 cursos participantes apenas 4 foram avaliados, cada um recebendo os conceitos 2, 3, 4 e máximo.
Nenhum curso recebeu o conceito mínimo e 2 ( 33,3 \% ) ficaram sem conceito. 5.2 Conceitos por Categoria Administrativa e por Região A Tabela 5.2 apresenta a distribuição dos cursos participantes do ENADE/2009 de DESIGN, por Categoria Administrativa, de acordo com os conceitos por eles alcançados, segundo a Região.
O Gráfico 5.2 apresenta a mesma informação sob a forma de diagrama de caixa para uma melhor visualização da dispersão dos conceitos obtidos.
Somente eventuais valores atípicos aparecem com o identificador da Região na qual o curso foi habilitado.
De acordo com as informações apresentadas, a grande maioria, 77,3 \%, é de cursos em instituições privadas, 12,7 \% dos cursos eram vinculados a instituições federais, 6,0 \% a instituições estaduais e os restantes 4,0 \% a instituições municipais.
As instituições estaduais foram as que apresentaram, em termos nacionais, maior proporção de cursos com conceito 5, conceito modal para esta categoria, equivalentes a 44,4 \% do total.
O conceito modal nas instituições federais foi 4, com 47,4 \% dos cursos.
Nas instituições municipais, o conceito modal foi 3, obtido por 50,0 \% dos cursos desta categoria.
Entre as instituições privadas o conceito modal foi 3, para 33,6 \% dos cursos.
Apenas as instituições privadas obtiveram conceito mínimo, equivalentes a 9,5 \% dos cursos desta categoria.
Ficaram sem conceito um total de 29 cursos, 21 de instituições privadas ( 18,1 \% da categoria ), 5 de instituições federais ( 26,3 \% da categoria ), 2 de instituições estaduais ( 22,2 \% da categoria ) e 1 de instituição municipal ( 16,7 \% da categoria ).
Considerando -se a análise por Região observou -se que as instituições privadas constituem a maioria em todas as Regiões Brasileiras.
Na Região Norte, dos 8 cursos participantes, 6 ( 75,0 \% ) eram de instituições privadas, 1 de instituição federal e 1 de instituição estadual, os dois últimos ficaram com conceito 4.
O conceito modal para cursos em instituições privadas foi o conceito 3 ( 50,0 \% ) e 1 curso obteve o conceito mínimo ( 16,7 \% ).
Na Região Nordeste, dos 15 cursos participantes, 8 ( 53,3 \% ) eram de instituições privadas, 6 de instituições federais ( 40,0 \% ) e 1 de instituição estadual ( 6,7 \% ).
Este último curso ficou com conceito 5.
Dos 6 cursos em instituições federais, 2 ficaram sem conceito e 4 receberam conceito 4.
O conceito modal para cursos em instituições privadas foi o conceito 3 ( 37,5 \% ) e 2 cursos obtiveram o conceito mínimo.
Também não foram avaliados na Região 2 cursos pertencentes a instituições federais.
Na Região Sudeste, dos 69 cursos participantes, 60 ( 87,0 \% ) eram de instituições privadas, 4 de instituições federais, 4 de instituições estaduais e 1 de instituição municipal.
O único curso em instituição municipal, recebeu conceito 3.
Dos 4 cursos em instituições estaduais, 2 ficaram sem conceito e 2 com conceito 4.
Dos 4 cursos em instituições federais, 1 ficou sem conceito, 1 recebeu o conceito 3 e 2 foram alocados na categoria modal, conceito 4.
O conceito modal para cursos em instituições privadas foi o conceito 3 ( 41,7 \% ) e 6 cursos ( 10,0 \% ) obtiveram o conceito mínimo.
Também não foram avaliados na Região 9 cursos pertencentes a instituições privadas.
Na Região Sul, dos 52 cursos participantes, 39 ( 75,0 \% ) eram de instituições privadas, 6 de instituições federais ( 11,5 \% ), 3 de instituições estaduais ( 5,8 \% ) e 4 de instituições municipais ( 7,7 \% ).
Dos 4 cursos em instituições municipais, a metade recebeu o conceito 3 e a outra metade conceito 4.
Os 3 cursos em instituições estaduais receberam conceito máximo.
Dos 6 cursos em instituições federais, 2 ficaram sem conceito e os 4 restantes ficaram divididos igualmente nos conceitos 4 e 5.
O conceito modal para cursos em instituições privadas foi o conceito 3 ( 41,7 \% ), 6 cursos ( 10,0 \% ) obtiveram o conceito mínimo e 1 curso o conceito máximo.
Também não foram avaliados na Região 9 cursos pertencentes a instituições privadas, 1 em federal e 2 em estadual.
Na Região Centro-Oeste, dos 6 cursos participantes, 3 ( 50,0 \% ) eram de instituições privadas, 2 de instituições federais e 1 de instituição municipal.
Este curso em instituição municipal não foi avaliado.
Os 2 cursos em instituições federais ficaram divididos entre os conceitos 3 e 5.
Dos 3 cursos em instituições privadas, 1 ficou sem conceito e os 2 restantes ficaram com os conceitos 2 e 4.
Nenhum curso, nesta Região, recebeu o conceito mínimo. 5.3 Conceitos por Organização Acadêmica e por Região Na Tabela 5.3 encontra -se a distribuição dos conceitos atribuídos aos cursos participantes do ENADE/2009 de DESIGN, por Organização Acadêmica, segundo as Regiões.
O Gráfico 5.3 apresenta a mesma informação sob a forma de diagrama de caixa para uma melhor visualização da dispersão dos conceitos obtidos.
Somente eventuais valores atípicos apareceram com o identificador da Categoria Administrativa.
Os cursos desta área estão principalmente alocados em Universidades ( 50,7 \% ), depois em Faculdades ( 32,7 \% ) e Centros Universitários ( 16,7 \% ).
Este ordenamento se repete em todas as Regiões, com exceção da Região Norte, na qual Universidades estão em segundo lugar.
Nenhum curso de Instituto Federal de Educação, Ciência e Tecnologia ( IFECT ) nem de CEFET participou.
Nas Universidades, o conceito modal foi 4 ( 30,3 \% ), nos Centros Universitários e Faculdades, o conceito modal foi 3 ( respectivamente 28,0 \% e 36,7 \% ).
Proporcionalmente poucos cursos nos 3 tipos de Organização Acadêmica mais frequentes, Universidades, Centros Universitários e Faculdades, receberam conceito mínimo, respectivamente 3,9 \%, 8,0 \% e 12,2 \%.
O conceito máximo foi alcançado por 10 cursos, 9 em Universidades e 1 em Faculdade.
Nos 3 tipos de Organização Acadêmica mais frequentes uma porcentagem considerável de cursos ficaram sem conceito, respectivamente 15,8 \%, 28,0 \% e 20,4 \%.
Considerando -se separadamente as cinco Regiões, observa -se que, a Região Norte dos 8 cursos participantes, 3 eram de Universidades e 5 de Faculdades.
Dos 3 cursos em Universidades, 2 receberam conceito 4 e 1 o conceito 3.
Todos os 5 cursos em Faculdades foram avaliados e os conceitos modais foram 2 e 3 ( 2 cursos cada ).
Como já comentado, nenhum curso nesta Região obteve conceito máximo.
O conceito mínimo foi alocado somente a Faculdades ( 1 curso ).
Na Região Nordeste dos 15 cursos participantes, 9 eram de Universidades, 1 de Centro Universitário e 5 de Faculdades.
O único curso de Centro Universitário recebeu conceito 3.
Dos 9 cursos em Universidades, 2 não foram avaliados, 5 receberam conceito 4, 1 o conceito máximo e 1 o conceito 2.
Dos 5 cursos em Faculdades, 2 receberam conceito 3, 1 o conceito 2 e 2 o conceito mínimo.
Na Região Sudeste, dos 69 cursos participantes, 29 eram de Universidades, 15 de Centros Universitários e 25 de Faculdades.
Dos 15 cursos de Centros Universitários, 3 não foram avaliados, os conceitos modais foram 2 e 3 ( 5 cursos cada ) e o conceito mínimo foi alocado a 2 cursos.
Dos 29 cursos em Universidades, 5 não foram avaliados e, entre os avaliados, o conceito modal foi 3 ( 10 cursos ).
Os 14 restantes receberam conceitos 1 ( 2 cursos ), 2 ( 6 cursos ) e 4 ( 6 cursos ).
Dos 25 cursos em Faculdades, 4 não foram avaliados.
Entre os cursos avaliados, o conceito modal foi também 3 ( 12 cursos ), 2 cursos receberam conceito mínimo, 2 conceito 2, 4 conceito 4 e 1 conceito máximo.
Na Região Sul dos 52 cursos participantes, 30 eram de Universidades, 9 de Centros Universitários e 13 de Faculdades.
Dos 9 cursos de Centros Universitários, 4 não foram avaliados, o conceito modal entre os avaliados foi 4 ( 3 cursos ), 1 curso recebeu conceito 2 e 1 curso o conceito 3.
Dos 30 cursos em Universidades, 3 não foram avaliados e, entre os avaliados, o conceito modal foi 4 ( 10 cursos ).
Os 17 restantes receberam conceitos 1 ( 1 curso ), 2 ( 2 cursos ), 3 ( 7 cursos ) e conceito máximo ( 7 cursos ).
Dos 13 cursos em Faculdades, 6 não foram avaliados.
Entre os cursos avaliados, o conceito modal foi também 4 ( 3 cursos ).
Nenhum curso em Faculdade recebeu conceito máximo nesta Região.
Na Região Centro-Oeste dos 6 cursos participantes, 5 eram de Universidades e 1 de Faculdade.
Dos 5 cursos em Universidades, 2 não foram avaliados e os 3 avaliados ficaram distribuídos nos conceitos 2, 3 e 5.
O único curso em Faculdade recebeu conceito 4.
Nenhum curso recebeu conceito mínimo nesta Região.

% matched lemmas: 
\end{textsample}
